
% =====================================================
% CONFIGURACIONES DE PAQUETES Y AJUSTES DEL DOCUMENTO
% =====================================================

% -----------------------------
% Idioma y codificación
% -----------------------------
% --- Configuración para babel ---
% Puedes cambiar el idioma principal o agregar otros
\selectlanguage{spanish} % Selecciona español como idioma principal

% -----------------------------
% Márgenes y geometría
% -----------------------------
% --- Configuración para geometry ---
% Define los márgenes y dimensiones de la página
\geometry{
    top=2.5cm,            % Margen superior
    bottom=2cm,           % Margen inferior
    left=2.5cm,           % Margen izquierdo
    right=2cm,            % Margen derecho
    heightrounded,        % Redondea la altura útil de la página
    bindingoffset=0.5cm   % Espacio extra para encuadernación
}

% -----------------------------
% Encabezados y pies de página
% -----------------------------
% --- Configuración para fancyhdr ---
% Personaliza el estilo de los encabezados y pies usando fancyhdr
\pagestyle{fancy}           % Usa el estilo 'fancy' para encabezados/pies
\fancyhf{}                  % Limpia encabezados y pies por defecto
\fancyhead[L]{ENCABEZADO}   % Texto a la izquierda del encabezado
\fancyhead[R]{\thepage}     % Texto a la derecha del encabezado
\fancyfoot[C]{\leftmark}    % Texto al centro del pie
\setlength{\headheight}{14.5pt} % Altura del encabezado (evita warnings)

% -----------------------------
% Títulos de sección y subsección
% -----------------------------
% --- Configuración para titlesec ---
% Modifica el formato de los títulos usando titlesec
	\titleformat{\section}
    {\bfseries\large}      % Título de sección en negrita y grande
    {\thesection. }        % Número de sección seguido de punto
    {0pt}{}                % Espacio entre número y título
	\titleformat{\subsection}
    {\bfseries}            % Título de subsección en negrita
    {\thesubsection. }     % Número de subsección seguido de punto
    {0pt}{}                % Espacio entre número y título

% -----------------------------
% Hipervínculos y PDF
% -----------------------------
% --- Configuración para hyperref ---
% Habilita enlaces en el PDF, define sus colores y propiedades
\hypersetup{
    colorlinks=true,      % Activa el color en los enlaces en vez de recuadros
    linkcolor=magenta,    % Color de los enlaces internos (índice, referencias)
    filecolor=magenta,    % Color de los enlaces a archivos
    urlcolor=magenta,     % Color de los enlaces a URLs externas
    pdftitle={Documento}, % Título del documento en las propiedades del PDF
    pdfpagemode=UseOutlines, % Abre el PDF mostrando el panel de marcadores
}

% -----------------------------
% Imágenes y figuras
% -----------------------------
% --- Configuración para graphicx ---
% Carpeta predeterminada para imágenes
\graphicspath{{figures/}} % Busca imágenes en la carpeta 'figures'

% --- Configuración para caption ---
% Personaliza el formato de los títulos de figuras y tablas
\captionsetup{
    font=small,             % Tamaño de fuente para captions
    labelfont=bf,           % Etiqueta (Figura, Tabla) en negrita
    justification=centering % Centra el texto del caption
}

% --- Configuración para float ---
% Permite forzar la posición exacta de figuras/tablas con [H]
% No requiere configuración adicional, pero puedes definir estilos si lo deseas

% -----------------------------
% Tablas
% -----------------------------
% (Puedes agregar configuraciones para booktabs, array, multirow, multicol si lo necesitas)

% -----------------------------
% Código fuente y pseudocódigo
% -----------------------------
% --- Configuración para listings (código fuente) ---
% Personaliza el estilo de los bloques de código
\lstset{
    basicstyle=\ttfamily\small,     % Fuente monoespaciada y tamaño pequeño
    numbers=left,                   % Números de línea a la izquierda
    numberstyle=\tiny,              % Estilo de los números de línea
    frame=single,                   % Marco alrededor del código
    breaklines=true,                % Rompe líneas largas automáticamente
    backgroundcolor=\color{gray!10},    % Fondo gris claro
    keywordstyle=\color{blue}\bfseries, % Palabras clave en azul y negrita
    commentstyle=\color{green!60!black}\itshape,    % Comentarios en verde y cursiva
    stringstyle=\color{red},        % Cadenas en rojo
}

% --- Configuración para algorithm2e (pseudocódigo) ---
% Personaliza el estilo de los algoritmos
\SetKwInput{KwData}{Datos}      % Renombra 'Input' a 'Datos'
\SetKwInput{KwResult}{Resultado} % Renombra 'Output' a 'Resultado'
\SetNlSty{textbf}{(}{)}         % Estilo para los números de línea

% -----------------------------
% Recuadros y colores
% -----------------------------
% --- Configuración para tcolorbox ---
% Define estilos para recuadros de colores
\tcbset{
    colback=blue!5!white,        % Fondo azul claro
    colframe=blue!80!black,      % Borde azul oscuro
    boxrule=0.8pt,               % Grosor del borde
    arc=4pt,                     % Bordes redondeados
    auto outer arc,              % Redondea también el borde exterior
}

% --- Configuración para xcolor ---
% Define colores personalizados
\definecolor{myred}{RGB}{220, 50, 47}    % Rojo personalizado
\definecolor{myblue}{RGB}{38, 139, 210}  % Azul personalizado

% -----------------------------
% Interlineado y espaciado
% -----------------------------
% --- Configuración para setspace ---
% Controla el interlineado del documento
% \singlespacing     % Interlineado sencillo
% \onehalfspacing    % Interlineado 1.5
% \doublespacing     % Interlineado doble

% --- Configuración de espaciado general ---
\setlength{\parskip}{1em}           % Espacio entre párrafos
\setlength{\parindent}{1.5em}       % Sangría al inicio de cada párrafo
\setlength\emergencystretch{3em}    % Estira líneas para evitar desbordes

% -----------------------------
% Notas al pie
% -----------------------------
% --- Configuración para footmisc ---
% Personaliza las notas al pie
\renewcommand{\thefootnote}{\fnsymbol{footnote}} % Usa símbolos en vez de números

% -----------------------------
% Bibliografía
% -----------------------------
% --- Configuración para biblatex ---
% Archivo de bibliografía y opciones adicionales
\addbibresource{referencias.bib}    % Archivo de bibliografía para biblatex
% \ExecuteBibliographyOptions{sorting=none, maxnames=4} % Ejemplo: sin orden y máximo 4 autores
